\bumppoints{26}

\question
Retinitis pigmentosus (RP) is a rare form of blindness, but has a high frequency on Tristen de Cu\~nha, a small island in the South Atlantic Ocean.  The island was colonized by 15 Britons during the 1800s.  One person carried one allele for RP.  Today, of the 240 people descended from the original 15, 4 are blind and 9 are carriers for the disease.  Which of the following best explains the high proportion of RP on the island?

\begin{choices}
\choice founder effect and natural selection.
\correctchoice founder effect and inbreeding (a form of non-random mating).
\choice inbreeding and mutation.
\choice genetic drift and mutation.
\choice bottleneck and inbreeding.
\end{choices}

\question
Which of these conditions should completely prevent the occurrence of natural selection in a population over time?

\begin{choices}
\choice The population size is large. 
\correctchoice All variation between individuals is due only to environmental factors. 
\choice The environment is changing at a relatively slow rate. 
\choice The population lives in a habitat where there are no competing species present. 
\choice The population mates non-randomly.
\end{choices}

\question
\textit{Anopheles} mosquitoes, which carry the malaria parasite, cannot live above elevations of 5,900 feet. In addition, oxygen availability decreases with higher altitude. Consider a hypothetical human population that is adapted to life on the slopes of Mt. Kilimanjaro in Tanzania, a country in equatorial Africa. Mt. Kilimanjaro's base is about 2,600 feet above sea level and its peak is 19,341 feet above sea level. If the frequency of the sickle-cell allele in the population is plotted against altitude (feet above sea level), which of the following distributions is most likely, assuming little migration of people up or down the mountain?

\begin{multicols}{2}
\begin{choices}
\choice \includegraphics[width=0.35\textwidth]{anopA} 
\correctchoice \includegraphics[width=0.35\textwidth]{anopB} 
\choice \includegraphics[width=0.35\textwidth]{anopD} 
\choice \includegraphics[width=0.35\textwidth]{anopC} 
\end{choices}
\end{multicols}

\question
According to Darwin's concept of descent with modification:

\begin{choices}
\choice Populations evolve by natural selection acting on individuals. 
\choice Individuals with greater fitness will evolve more quickly than individuals with lower fitness. 
\correctchoice Species evolve from ancestral forms through the accumulation of different adaptations. 
\choice All organisms overproduce offspring, leading to different adaptations. 
\choice All organisms overproduce offspring, leading to differences in survivorship. 
\end{choices}

\question
Snails living in the southern rocky intertidal have adapted to warm temperatures. Compared to northern snails, southern snails have ridges to get rid of excess heat and lighter colors to reflect heat. If the snail phenotype ranges from no ridges and dark color to ridges and light color, then the mode of selection that fits the southern snails is most likely

\begin{choices}
\choice disruptive selection. 
\choice sexual selection. 
\correctchoice directional selection. 
\choice stabilizing selection. 
\choice balancing selection. 
\end{choices}

\question
Violation of which two Hardy-Weinberg assumptions \emph{always} causes overall homozygosity to increase in a population?

\begin{choices}
\choice Natural selection and gene flow.
\correctchoice Genetic drift and non-random mating.
\choice Mutation and natural selection.
\choice Non-random mating and gene flow.
\choice Genetic drift and natural selection.
\choice Gene flow and mutation.
\end{choices}

\question
Which mouse has the greatest relative fitness? (A clutch is the offspring from one breeding event.)

\begin{choices}
\choice One that lives for 4 years, has 3 clutches per year, and 4 offspring per clutch.
\choice One that lives for 3 years, has 2 clutches per year, and 6 offspring per clutch. 
\correctchoice One that lives for 3 years, has 3 clutches per year, and 7 offspring per clutch. 
\choice One that lives for 2 years, has 3 clutches per year, and 8 offspring per clutch. 
\choice One that lives for 3 years, has 3 clutches per year, and 5 offspring per clutch.
\end{choices}

\question
The males of some fishes and birds have bright colors that males use during territorial displays. The existence of these bright colors in unrelated organisms is an example of 

\begin{choices}
\correctchoice convergent evolution. 
\choice common ancestry. 
\choice intersexual selection. 
\choice descent with modification. 
\choice homology. 
\end{choices}

\question
Prairie chickens are endangered grassland birds. In Illinois, 
grasslands used to cover most of the state but farming and growth of 
cities changed the habitat over a fairly short period of time. Now, 
less than 1\% of the grasslands remain. A study showed that most genes 
had far fewer alleles than historic populations. The loss of genetic 
variation is most likely due to

\begin{choices}
\choice a founder event caused by individuals moving to the remaining grasslands. 
\choice the small population size causing random fluctuation of allele frequencies. 
\choice balancing selection, which maintains genetic variation in small populations. 
\correctchoice a genetic bottleneck caused by the loss of habitat. 
\choice directional selection that favored survivorship in small habitats. 
\end{choices}

\question
The overuse of drugs such as penicillin has caused an increase in the number of bacterial species that are resistant to the drugs. The most likely explanation for this result is:

\begin{choices}
\choice The drugs caused mutations in the bacteria that gave them resistance. 
\choice Most bacteria can resist drugs. 
\choice Bacteria began making proteins resistant to the drugs. 
\correctchoice A few bacteria were naturally resistant and natural selection favored their increase. 
\choice Penicillin was not an effective drug. 
\end{choices}

\question
The rabbitsfoot mussel was once widespread throughout rivers of eastern North America but now remains in only a few very small populations. The decrease occurred rapidly due to habitat loss. Genetics studies show that the remaining populations have little or no genetic variation. The best explanation for this result is:

\begin{choices}
\choice Inbreeding caused an increase in heterozygosity. 
\choice A bottleneck caused the loss of homozygosity. 
\choice Balancing selection favored only one or two alleles. 
\correctchoice Genetic drift caused the loss of heterozygosity. 
\choice Directional selection favored only one or two alleles. 
\end{choices}

\question
A species of sport fish reaches an average adult length of 100 cm. Fishing regulations require that only individuals longer than 75 cm may be kept. Shorter fish must be released because they have greater reproductive output relative to large individuals. Based on your knowledge of natural selection, you should predict that over hundreds of generations the average length of the adult fish population  will

\begin{choices}
	\choice gradually increase. 
	\choice gradually decline. 
	\correctchoice remain unchanged. 
	\choice rapidly decline. 
	\choice rapidly increase. 
\end{choices}

\newpage

\question
Which \emph{one} of the following statements is true of natural selection?

\begin{choices}
	\choice It requires genetic variation. 
	\choice It results in descent with modification. 
	\choice It involves differences in reproductive success. 
	\correctchoice It requires genetic variation, results in descent with modification, and involves differences in reproductive success. 
	\choice It results in descent with modification and involves differences in reproductive success. 
\end{choices}
